\section{Conclusion}\label{conclusion}

This paper studies upper-tail food-price forecasting when the
environment changes and forecast errors arrive with a three-month delay.
Static q90 forecasts fail severely during Nigeria's 2023 food-price
regime change. Delayed EWMA and conformal residual calibration recover
part of the lost calibration, but no single adaptive method is best
across sources and regimes.

A regime-aware controller that combines static, gated, global conformal,
and hierarchical experts has lower pooled post-2022 loss than every
individual expert in the NBS continuation and WFP external panel. It
significantly outperforms static forecasting, but not the best expert,
and it overforecasts during the NBS 2025 recovery. The defensible
conclusion is therefore that expert aggregation can reduce the cost of
choosing one calibration rule under delayed feedback, while recovery and
subgroup calibration remain unresolved.

The next scientific step is not further tuning on the Nigerian data. It
is a preregistered, frozen evaluation on an untouched country or future
period. Until that test is completed, the controller should be presented
as a promising exploratory forecasting system rather than a confirmed
universally superior method.
